% Part: normal-modal-logic
% Chapter: frame-correspondence
% Section: properties-accessibility

\documentclass[../../../include/open-logic-section]{subfiles}

\begin{document}

\olfileid[pt]{nml}{frd}{acc}

\olsection{Propriedades das relações de acessibilidade}

Muitas !!{formula}s modais revelam-se características de propriedades simples
e até familiares da relação de acessibilidade. Em uma direção, isso significa
que todo modelo que possui determinada propriedade torna verdadeira uma
!!{formula} correspondente (e todas as suas instâncias de substituição).
Começamos com cinco exemplos clássicos de tipos de relações de acessibilidade
e das fórmulas cuja verdade elas garantem.

\begin{thm}\ollabel{thm:soundschemas}
  Seja $\mModel{M} = \tuple{W, R, V}$ um modelo. Se $R$ possui a propriedade
  no lado esquerdo de \olref{tab:five}, toda instância da !!{formula} no lado
  direito é verdadeira em~$\mModel{M}$.
\end{thm}

\begin{table}[t]
    \begin{tabular}{| l || l |}
      \hline
      {\emph{Se $R$ é \dots}} & {\emph{então \dots é verdadeira em~$\mModel{M}$:}} \\
      \hline \hline
      \emph{serial}: $\forall u \exists v Ruv$ & \hbox to.43\textwidth
           {$\Box p \lif \Diamond p$ \hfill (\Ax{D})} \\
      \hline
      \emph{reflexiva}: $\forall w Rww$  
      & \hbox to.43\textwidth
          {$\Box p \lif p$ \hfill (\Ax{T})} \\
      \hline
      \emph{simétrica}: &  \hbox to .43\textwidth
           {$p \lif \Box\Diamond p$ \hfill (\Ax{B})} \\
           $\forall u\forall v(Ruv \lif Rvu)$ & \\
      \hline
      \emph{transitiva}: & \hbox to.43\textwidth
           {$\Box p \lif \Box \Box p$ \hfill (\Ax{4})} \\
      $\forall u \forall v \forall w ((Ruv \land Rvw) \lif Ruw)$ & \\
      \hline 
      \emph{euclidiana}: & \hbox to .43\textwidth
        {$\Diamond p \lif \Box\Diamond p$ \hfill (\Ax{5})} \\
      $\forall w \forall u \forall v ((Rwu \land Rwv) \lif Ruv)$ &\\
      \hline
    \end{tabular}
    \caption{Cinco fatos de correspondência.}
    \ollabel{tab:five}
  \end{table} 


\begin{proof}
  Eis o caso de \Ax{B}: para mostrar que o esquema é verdadeiro em um modelo,
  precisamos mostrar que todas as suas instâncias são verdadeiras em todos os
  mundos do modelo. Seja, então, $!A \lif \Box\Diamond !A$ uma instância dada
  de \Ax{B}, e seja $w \in W$ um mundo arbitrário. Suponha que o antecedente
  $!A$ seja verdadeiro em $w$, a fim de mostrar que $\Box \Diamond !A$ é
  verdadeira em $w$. Portanto, precisamos mostrar que $\Diamond !A$ é
  verdadeira em todos os $w'$ acessíveis a partir de $w$. Agora, para todo
  $w'$ tal que $Rww'$, temos, usando a hipótese de simetria, também $Rw'w$
  (veja \olref{fig:Bsymm}). Como $\mSat{M}{!A}[w]$, temos
  $\mSat{M}{\Diamond !A}[w']$. Como $w'$ era um mundo arbitrário tal que
  $Rww'$, temos $\mSat{M}{\Box\Diamond!A}[w]$.

  Deixamos os outros casos como exercícios.
\end{proof}

\begin{prob}
  Complete a demonstração de \olref[nml][frd][acc]{thm:soundschemas}.
\end{prob}

\begin{figure}
  \begin{center}
    \begin{tikzpicture}[modal]
      \node[world] (w1) [label={below:$\mSat{{}}{\formula{A}}$\\
          $\mSat{{}}{\Box\Diamond \formula{A}}$}] {$w$}; 
      \node[world] (w2) [label=below:{$\mSat{{}}{\Diamond \formula{A}}$},
        right=of w1] {$w'$}; 
      \draw[->, bend left] (w1) to (w2); 
      \draw[->, bend left] (w2) to (w1); 
    \end{tikzpicture}
  \end{center}
\caption{O argumento a partir da simetria.}
\ollabel{fig:Bsymm}
\end{figure}

Observe que as implicações recíprocas de \olref{thm:soundschemas} não valem:
não é verdade que, se um modelo verifica um esquema, então a relação de
acessibilidade desse modelo possui a propriedade correspondente. No caso de
\Ax{T} e dos modelos reflexivos, é fácil dar um exemplo de modelo no qual a
própria \Ax{T} falha: sejam $W = \{w\}$ e $V(p) = \emptyset$. Então $R$ não é
reflexiva, mas $\mSat{M}{\Box p}[w]$ e $\mSat/{M}{p}[w]$. Aqui, porém, temos
apenas uma instância de \Ax{T} que falha em~$\mModel{M}$; outras instâncias,
por exemplo, $\Box \lnot p \lif \lnot p$, são verdadeiras. É mais difícil
dar exemplos nos quais \emph{toda instância de substituição} de~\Ax{T} seja
verdadeira em~$\mModel{M}$ e $\mModel{M}$ não seja reflexivo. Mas também
existem tais modelos:

\begin{prop}\ollabel{prop:reflexive}
  Seja $\mModel{M} = \tuple{W, R, V}$ um modelo tal que $W = \{u, v\}$, no
  qual os mundos $u$ e $v$ estão relacionados por $R$: isto é, tanto $Ruv$
  como $Rvu$. Suponha que, para todo $p$: $u \in V(p) \Leftrightarrow v
  \in V(p)$. Então:
  \begin{enumerate}
  \item Para todo $!A$: $\mSat{M}{!A}[u]$ se e somente se
    $\mSat{M}{!A}[v]$ (use indução sobre $!A$).
  \item Toda instância de~\Ax{T} é verdadeira em $\mModel{M}$.
  \end{enumerate}
  Como $\mModel{M}$ não é reflexivo (na verdade, ele é
  \emph{irreflexivo}), a recíproca de \olref{thm:soundschemas} falha no caso
  de \Ax{T} (argumentos semelhantes podem ser apresentados para alguns---mas
  não todos---os outros esquemas mencionados em \olref{thm:soundschemas}).
\end{prop}

\begin{prob}
  Demonstre as afirmações de \olref[nml][frd][acc]{prop:reflexive}.
\end{prob}

Embora nos concentremos nas cinco !!{formula}s clássicas \Ax{D}, \Ax{T},
\Ax{B}, \Ax{4} e \Ax{5}, registramos em \olref{tab:anotherfive} algumas
outras propriedades de relações de acessibilidade. A relação de
acessibilidade~$R$ é parcialmente funcional se, de cada mundo, no máximo um
mundo é acessível. Se, de cada mundo, exatamente um mundo é acessível,
chamamos a relação de funcional. (Assim, as relações funcionais são
precisamente aquelas que são ao mesmo tempo seriais e parcialmente
funcionais.) Elas são chamadas de «funcionais» porque a relação de
acessibilidade opera como uma função (parcial). Uma relação é fracamente
densa se, sempre que $Ruv$, existe um~$w$ «entre» $u$ e~$v$. Assim, as
relações fracamente densas são, em certo sentido, o oposto das relações
transitivas: em uma relação transitiva, sempre que se pode chegar a $v$ a
partir de~$u$ por um desvio via~$w$, pode-se chegar a $v$ a partir de $u$
diretamente; em uma relação fracamente densa, sempre que se pode chegar a $v$
a partir de~$u$ diretamente, também se pode chegar a ele por um desvio via
algum~$w$. Uma relação é fracamente dirigida se, sempre que se pode chegar aos
mundos $u$ e~$v$ a partir de algum mundo~$w$, pode-se chegar a um único
mundo~$t$ a partir tanto de $u$ como de~$v$---isso às vezes é chamado de
«propriedade do losango» ou «confluência».

\begin{table}[t]
    \begin{tabular}{| p{.50\textwidth} || p{.43\textwidth} |}
      \hline
      {\emph{Se $R$ é \dots}} & {\emph{então \dots é verdadeira em~$\mModel{M}$:}} \\
      \hline \hline
      \emph{parcialmente funcional}: & 
      \multirow{2}{*}{$\Diamond p \lif \Box p$} \\
      $ \forall w \forall u \forall v ((Rwu \land Rwv) \lif u =v )$ & \\
      \hline
      \multirow{2}{*}{\emph{funcional}: $\forall w \exists v \forall u(Rwu \liff \eq[u][v]) $} 
      & \multirow{2}{*}{$\Diamond p \liff \Box p$} \\
      & \\
      \hline
      \emph{fracamente densa}: &  \multirow{2}{*}{$\Box \Box p \lif
        \Box p$} \\
      $\forall u\forall v(Ruv \lif \exists w(Ruw \land Rwv))$ & \\
      \hline
      \emph{fracamente conexa}: &
      \multirow{3}{*}{\hbox to.43\textwidth{$
        \begin{array}{@{}l@{}}
          \Box ((p \land \Box p) \lif q) \lor {} \\
          \qquad \Box ((q \land \Box q) \lif p)
        \end{array}$ \hfill (\Ax{L})}
      } \\
      $\forall w \forall u \forall v ((Rwu \land Rwv) \lif {}$ &\\
      \qquad $(Ruv \lor u=v \lor Rvu))$ & \\
      \hline 
      \emph{fracamente dirigida}: & \multirow{3}{*}{\hbox to .43\textwidth
        {$\Diamond\Box p \lif \Box\Diamond p$\hfill (\Ax{G})}} \\
      $\forall w \forall u \forall v ((Rwu \land Rwv) \lif {}$ & \\
      \qquad $\exists t (Rut \land Rvt))$ & \\
      \hline
    \end{tabular}
    \caption{Mais cinco fatos de correspondência.}
    \ollabel{tab:anotherfive}
  \end{table} 

\begin{prob}
  Seja $\mModel{M} = \tuple{W, R, V}$ um modelo. Mostre que, se $R$ satisfaz
  as propriedades do lado esquerdo de
  \olref[nml][frd][acc]{tab:anotherfive}, toda instância da fórmula
  correspondente do lado direito é verdadeira em~$\mModel{M}$.
\end{prob}



\end{document}
